\documentclass[11pt]{amsart}
\usepackage[margin=1.15in]{geometry}
\usepackage[T1]{fontenc}
\usepackage{lmodern,microtype}
\usepackage{amsmath,amssymb,amsthm,mathtools,booktabs}
\usepackage[colorlinks=true,linkcolor=blue,citecolor=blue,urlcolor=blue]{hyperref}

\newtheorem{theorem}{Theorem}[section]
\newtheorem{lemma}[theorem]{Lemma}
\newtheorem{proposition}[theorem]{Proposition}
\newtheorem{corollary}[theorem]{Corollary}
\newtheorem{conjecture}[theorem]{Conjecture}
\theoremstyle{remark}
\newtheorem{remark}[theorem]{Remark}

\newcommand{\Z}{\mathbb{Z}}
\newcommand{\R}{\mathbb{R}}
\newcommand{\Res}{\operatorname{Res}}

\title{Ehrhart positivity of base-$r$ simplices}
\author{Brandon Li}
\address{Great Neck, New York, USA}
\email{brandon.li.gn@gmail.com}
\subjclass[2020]{52B20, 05A15, 11B68}
\keywords{Ehrhart polynomial, Ehrhart positivity, lattice simplex, numeral system, generating function}

\begin{document}
\begin{abstract}
For an integer $r\ge2$, the base-$r$ simplex $B_{r,d}$ is the convex hull of
the standard basis vectors of $\R^d$ and the point $-(r-1)(1,r,\ldots,r^{d-1})$.
A lattice polytope is Ehrhart positive if every coefficient of its Ehrhart
polynomial is positive. Solus conjectured that every $B_{r,d}$ is Ehrhart
positive, and Graham showed that $B_{14,5}$ is not. We prove that $B_{r,d}$ is
Ehrhart positive in every dimension if and only if $r\in\{2,3,4,5\}$, and for
each $r\ge6$ we find the smallest dimension in which a negative coefficient
appears. We also determine the eventual sign of every coefficient when the
dimension is fixed and the base is large, locate the exact sign change of the
coefficient four degrees below the leading term at $r-1=(5d-3)/2$, and show
that at any fixed base the coefficients at fixed distance below the leading
degree are eventually positive. The proofs combine a generating function for
the Ehrhart polynomials, a residue formula over the $(r-1)$st roots of unity,
and, for bases four and five, Lagrange inversion with explicit analytic tail
bounds.
\end{abstract}
\maketitle

\section{Introduction}

For integers $r\ge2$ and $d\ge1$, set $q=r-1$ and define the $d$-dimensional
lattice simplex
\[
 B_{r,d}=\operatorname{conv}\bigl(e_1,\ldots,e_d,\;-q(1,r,r^2,\ldots,r^{d-1})\bigr)\subset\R^d .
\]
Write $L_d(t)=\#\bigl(tB_{r,d}\cap\Z^d\bigr)$ for its Ehrhart polynomial; the
base $r$ will always be clear from context. Solus~\cite{Sol} introduced these
simplices in connection with base-$r$ numeral systems, computed their
$h^*$-polynomials (the numerator polynomials of their Ehrhart series), and
proved that those polynomials have only real roots. Based on numerical
evidence, he also suggested that these simplices are Ehrhart positive, meaning
that every coefficient of $L_d(t)$ is positive. Liu recorded this as a
conjecture of Solus in her survey~\cite{Liu}, and Liu and
Solus~\cite[Conjecture~2.2]{LS} stated it for all $r\ge2$ and $d\ge1$.

Graham~\cite{Gra} disproved the conjecture by exhibiting a negative linear
coefficient:
\[
 L_5(t)=1-\tfrac{268}{15}t+\tfrac{3955}{12}t^2+\tfrac{5833}{2}t^3
 +\tfrac{10343}{12}t^4+\tfrac{67228}{15}t^5\qquad(r=14).
\]
This leaves two natural questions: which bases are Ehrhart positive in every
dimension, and how are the negative coefficients distributed when positivity
fails? Our first main result completely answers the first question.

\begin{theorem}\label{thm:classification}
The simplex $B_{r,d}$ is Ehrhart positive for every $d\ge1$ if and only if
$r\in\{2,3,4,5\}$. For $r\ge6$ the least dimension $d$ in which $L_d(t)$ has a
negative coefficient is given in Table~\ref{tab:first}; in that dimension the
negative coefficient is the linear one.
\end{theorem}

\begin{table}[ht]
\centering
\begin{tabular}{c|ccccc}
\toprule
$r$ & $6$ & $7$ & $8$ & $9,\ldots,13$ & $\ge14$\\
\midrule
first $d$ with a negative coefficient & $43$ & $19$ & $15$ & $7$ & $5$\\
\bottomrule
\end{tabular}
\caption{The first failure of Ehrhart positivity for $B_{r,d}$, $r\ge6$.}
\label{tab:first}
\end{table}

The case $r=2$ was already known. Solus~\cite{Sol} showed that $B_{2,d}$ has
$h^*$-polynomial $(1+z)^d$, which coincides with that of the $d$-dimensional
cross-polytope. The two polytopes therefore share the same Ehrhart polynomial,
whose zeros lie on the line $\operatorname{Re}t=-1/2$~\cite{BCKV}, implying
Ehrhart positivity. Section~\ref{sec:23} gives a short, direct proof for both
$r=2$ and $r=3$. The main effort of the paper is the proof for $r=4$ and $r=5$
in Section~\ref{sec:45}.

The sign of the coefficient of $t^k$ is governed by its codimension $d-k$, the
difference between the degree and the exponent. When the base is large relative
to the dimension, a periodic sign pattern emerges.

\begin{theorem}\label{thm:large}
Fix $d\ge1$. For all sufficiently large $r$, the coefficient of $t^k$ in
$L_d(t)$, $0\le k\le d$, is negative if and only if
\[
 k\ge1,\qquad d-k\ge4,\qquad d-k\equiv0\text{ or }1\pmod4,
\]
and is positive otherwise.
\end{theorem}

Negative coefficients first appear in codimension four. In this codimension, we
can determine the exact threshold where the sign changes.

\begin{theorem}\label{thm:codim4}
Let $q=r-1$. For $d\ge6$, the coefficient of $t^{d-4}$ in $L_d(t)$ is negative
if $q>(5d-3)/2$ and positive if $q\le(5d-3)/2$. For $d=5$ it is negative if and
only if $r\ge14$. For $d\le4$ and every $r\ge2$, all coefficients of $L_d(t)$
are positive.
\end{theorem}

When the base is held fixed, however, negative coefficients eventually
disappear in any fixed codimension. To state this behavior and the intermediate
scaling regime where both the base and dimension grow together, we rescale the
coefficients. For $0\le\ell\le d$, set
\[
 A_{d,\ell}=(d-\ell)!\,[t^{d-\ell}]\,L_d(t).
\]
Here, as throughout the paper, square brackets extract the indicated
coefficient of a polynomial or power series.

\begin{theorem}\label{thm:fixed}
Fix $\ell\ge0$.
\begin{enumerate}
\item If $r\to\infty$ and $d\to\infty$ with $r/d\to c\in(0,\infty)$, then
\[
 r^{-d}A_{d,\ell}\longrightarrow[v^\ell]\exp\Bigl(\frac{(v/2)\coth(v/2)-1}{c}\Bigr).
\]
For odd $\ell$ the limit is $0$, and $r^{1-d}A_{d,\ell}\to\frac12[v^{\ell-1}]\exp\bigl(((v/2)\coth(v/2)-1)/c\bigr)$.
\item For each fixed $r\ge2$, $A_{d,\ell}>0$ for all sufficiently large $d$.
\end{enumerate}
\end{theorem}

Expanding $(v/2)\coth(v/2)=1+v^2/12-v^4/720+O(v^6)$ shows that the coefficient
of $v^4$ in the limit in~(1) is $(5-2c)/(1440c^2)$. This expression changes
sign at $c=5/2$, matching the threshold in Theorem~\ref{thm:codim4}.

\subsection*{Outline}
In Section~\ref{sec:gf}, we express $L_d(t)$ as the coefficient of $x^{qt}$ in
an explicit rational function. Applying the residue theorem then rewrites the
rescaled coefficients as a sum over the $q$th roots of unity. From this setup
we derive three results: Theorem~\ref{thm:large} follows from the leading
terms in $q$ (Section~\ref{sec:large}), Theorem~\ref{thm:codim4} from an exact
formula for $A_{d,4}$ (Section~\ref{sec:codim4}), and Theorem~\ref{thm:fixed}
from the dominant contribution of the principal root of unity, namely one
(Section~\ref{sec:fixed}).

For the bases $2,3,4,5$, we instead sum the generating function over the
dimension. This yields a finite sum $\sum_\zeta b_\zeta(z)e^{tv_\zeta(z)}$
over the branches of an algebraic function. For $r=2,3$, every coefficient of
the principal branch is positive and the other branch has nonnegative
coefficients (Section~\ref{sec:23}). For $r=4,5$, we prove that the principal
branch dominates (Section~\ref{sec:45}). The first $180$ coefficients (for base
four) or $900$ coefficients (for base five) of every branch are computed
exactly via Lagrange inversion, and the remaining coefficients are bounded by
Cauchy estimates on a disk of radius $9/8$ or $51/50$. Finally,
Section~\ref{sec:classification} combines these ingredients to prove
Theorem~\ref{thm:classification}.

\section{Generating functions}\label{sec:gf}

Set $S_j(x)=1+x+\cdots+x^j$. Solus~\cite{Sol} showed that the $h^*$-polynomial
of $B_{r,d}$ is obtained from $S_q(x)^d=\sum_e c_ex^e$ by replacing each
$x^e$ with $z^{\lceil e/q\rceil}$:
\[
 h^*(z)=\sum_{e\ge0}c_e\,z^{\lceil e/q\rceil}.
\]
We first translate this description into a rational generating function for the
Ehrhart polynomial.

\begin{lemma}\label{lem:gf}
For all integers $t\ge0$,
\begin{equation}\label{eq:gf}
 L_d(t)=[x^{qt}]\,F_d(x),\qquad
 F_d(x)=\frac{1}{1-x}\Bigl(\frac{S_q(x)}{1-x^q}\Bigr)^{d}
 =\frac1{1-x}\Bigl(\frac1{1-x}+\frac{x^q}{1-x^q}\Bigr)^{d}.
\end{equation}
\end{lemma}

\begin{proof}
By the definition of $h^*$, we have $L_d(t)=\sum_j h^*_j\binom{t-j+d}{d}$,
with $\binom{n}{d}=0$ for $0\le n<d$. Write each positive exponent $e>0$ as
$e=qj-s$ with $j=\lceil e/q\rceil$ and $0\le s<q$, and for $e=0$ set $j=s=0$.

Define $G(x)=S_{q-1}(x)(1-x^q)^{-d-1}$. Expanding this series gives
$[x^{qm+s}]G=\binom{m+d}{d}$ for all $m\ge0$ and $0\le s<q$. Hence
$[x^{qt-e}]G=\binom{t-j+d}{d}$ whenever $e\le qt$. When $e>qt$, both sides
vanish: we have $j>t$, and since $j\le d$, it follows that $0\le t-j+d<d$.
Summing over all exponents therefore gives $L_d(t)=[x^{qt}]S_q(x)^dG(x)$.
Using $S_{q-1}(x)/(1-x^q)=1/(1-x)$ yields the first form in~\eqref{eq:gf}, and
the second follows from $1/(1-x)+x^q/(1-x^q)=S_q(x)/(1-x^q)$.
\end{proof}

Expanding the $d$th power in the second form of~\eqref{eq:gf} by the binomial
theorem and extracting coefficients gives
\begin{equation}\label{eq:conv}
 L_d(t)=\binom{qt+d}{d}+\sum_{j=1}^{d}\binom dj\sum_{u=0}^{t-1}
 \binom{t-u-1}{j-1}\binom{qu+d-j}{d-j}.
\end{equation}
Here, for $n\ge0$, $\binom{n}{k}$ denotes the polynomial
$n(n-1)\cdots(n-k+1)/k!$ in $n$, which vanishes when $0\le n<k$. By
Faulhaber's formula for sums of powers, the right-hand side is a polynomial in
$(q,t)$ of degree at most $d$ in $q$. Consequently, each coefficient
$[t^k]L_d(t)$ is a polynomial in $q$ of degree at most $d$.

Next we express the rescaled coefficients as a sum over the $q$th roots of
unity. For each $q$th root of unity $\zeta$, set
\[
 a_\zeta(v)=\frac{1}{1-\zeta e^{-v/q}},\qquad b(v)=\frac{1}{e^v-1}.
\]

\begin{lemma}\label{lem:residue}
For $0\le\ell\le d$,
\[
 A_{d,\ell}=[v^\ell]\sum_{\zeta^q=1}\frac vq\,a_\zeta(v)\bigl(v(a_\zeta(v)+b(v))\bigr)^d .
\]
\end{lemma}

\begin{proof}
Fix an integer $t\ge0$. The coefficient $[x^{qt}]F_d$ equals the residue of
$F_d(x)x^{-qt-1}$ at $x=0$. The only other poles of this rational function are
at the points $x=\zeta$ with $\zeta^q=1$, together with a possible pole at
$\infty$. As $x\to\infty$, we have $F_d(x)=O(x^{-1})$, so the residue at
infinity vanishes. By the residue theorem,
$[x^{qt}]F_d=-\sum_\zeta\Res_{x=\zeta}$.

At each pole, substitute $x=\zeta e^{-v/q}$, so that $dx=-(x/q)\,dv$,
$1/(1-x)=a_\zeta(v)$, $x^q/(1-x^q)=b(v)$, and $x^{-qt}=e^{tv}$. This yields
\[
 -\Res_{x=\zeta}F_d(x)x^{-qt-1}
 =\frac1q[v^{-1}]\,a_\zeta(a_\zeta+b)^de^{tv}
 =[v^{d}]\,\frac vq a_\zeta\bigl(v(a_\zeta+b)\bigr)^de^{tv}.
\]
Summing this expression over all roots of unity produces a polynomial in $t$
that agrees with the Ehrhart polynomial at every integer $t\ge0$, so the two
polynomials are identical. In each summand, the coefficient of $t^{d-\ell}$ is
$[v^\ell]\frac vq a_\zeta(v(a_\zeta+b))^d/(d-\ell)!$, which gives the lemma.
\end{proof}

In Lemma~\ref{lem:residue}, the principal root $\zeta=1$ behaves differently
from the others. At this root, we have the expansions
\begin{equation}\label{eq:principal}
 \frac vq a_1(v)=1+\frac{v}{2q}+\frac{v^2}{12q^2}+\cdots,\qquad
 v\bigl(a_1(v)+b(v)\bigr)=(q+1)\Bigl(1+\frac{v^2}{12q}+O(v^4)\Bigr).
\end{equation}
At every other root, $v\,a_\zeta(v)=O(v)$ and $v(a_\zeta+b)=1+O(v)$. Thus the
principal term carries the exponential factor $(q+1)^d=r^d$, whereas the
remaining terms grow only polynomially in $d$.

\section{Large base: proof of Theorem~\ref{thm:large}}\label{sec:large}

As shown in Section~\ref{sec:gf}, each coefficient $[t^k]L_d$ is a polynomial
in $q$ of degree at most $d$. For large $q$, its sign is therefore the sign of
its leading coefficient, which we read off from~\eqref{eq:conv}. Write
$S_n(t)=\sum_{u=0}^{t-1}u^n$ for the power sum. By Faulhaber's formula,
\[
 (n+1)S_n(t)=\sum_{i=0}^{n}\binom{n+1}{i}B_it^{n+1-i},
\]
where the Bernoulli numbers begin with $B_0=1$, $B_1=-1/2$, $B_2=1/6$, and so
on. Recall that $B_i=0$ for every odd $i\ge3$, and that $B_{2a}$ has sign
$(-1)^{a+1}$ for all $a\ge1$.

The only contribution to $q^d$ comes from $j=0$ and equals $q^dt^d/d!$. The
contributions to $q^{d-1}$ come only from $j=0$ and $j=1$:
\[
 (d-1)!\,[q^{d-1}]L_d(t)=\frac{d+1}2t^{d-1}+dS_{d-1}(t)
 =t^d+\frac12t^{d-1}+\sum_{i=2}^{d-1}\binom diB_it^{d-i}.
\]
Thus the leading coefficient is positive for $k=d$ and $k=d-1$. For an even
codimension $d-k=i\ge2$ with $k\ge1$, the leading coefficient is
$\binom diB_i/(d-1)!$, which has the same sign $(-1)^{i/2+1}$ as $B_i$. It is
therefore positive when $i\equiv2$ and negative when $i\equiv0\pmod 4$.

For an odd codimension $d-k=i\ge3$, the coefficient of $q^{d-1}t^k$ vanishes,
so we pass to $q^{d-2}$. Set $n=d-2$. In~\eqref{eq:conv}, the term $j=0$
contributes only a multiple of $t^{d-2}$ to $q^{d-2}$. The term $j=1$
contributes $d\cdot\frac{d(d-1)}{2}S_n(t)/(d-1)!$, and $j=2$ contributes
$\binom d2\bigl((t-1)S_n(t)-S_{n+1}(t)\bigr)/n!$. Combining these terms gives
\[
 [q^{d-2}]L_d(t)=\frac{d}{2\,n!}\Bigl(S_n(t)+(d-1)\bigl(tS_n(t)-S_{n+1}(t)\bigr)\Bigr)
 +c\,t^{d-2}
\]
for some constant $c$. Now take $k=d-i$ with $i=2a+1\ge3$ and $k\ge1$. The
coefficient of $t^k$ in $S_n$ involves $B_{i-1}=B_{2a}$, whereas the
corresponding coefficients in $tS_n$ and $S_{n+1}$ involve $B_i=0$. Hence the
leading coefficient is $\frac{d}{2\,n!}\binom{d-1}{2a}\frac{B_{2a}}{d-1}$. Its
sign is $(-1)^{a+1}$, which is positive for $i\equiv3$ and negative for
$i\equiv1\pmod4$.

Finally, the constant term is $L_d(0)=1$. Collecting all cases, the
coefficients that are negative for large $q$ are precisely those with $k\ge1$
and $d-k\in\{4,5,8,9,12,13,\ldots\}$. \qed

\section{Codimension four: proof of Theorem~\ref{thm:codim4}}\label{sec:codim4}

To locate the exact sign change in codimension four, we first derive a
closed-form expression for the rescaled coefficient.

\begin{proposition}\label{prop:codim4}
For all $d\ge1$ and $q\ge1$,
\begin{equation}\label{eq:codim4}
 4320\,q^4A_{d,4}=6(1+q)^dP+(q-1)T,
\end{equation}
where
\begin{align*}
P={}&-1+4dq+\Bigl(d+5\binom d2\Bigr)q^2-dq^3,\\
T={}&d(d-1)(d-2)q^6+d(d-1)(d-11)q^5+(-19d^3+18d^2+19d)q^4\\
 &+(26d^3+18d^2+4d-6)q^3+(36d^2+12d-6)q^2+(18d-6)q-6.
\end{align*}
\end{proposition}

\begin{proof}
We expand the formula in Lemma~\ref{lem:residue} at $\ell=4$. By~\eqref{eq:principal},
the principal term $\zeta=1$ equals $(1+q)^d[v^4]\,U(1+w)^d$, where
$U=\frac vqa_1$ and $1+w=v(a_1+b)/(q+1)$. Expanding
$(1+w)^d=\exp(d\log(1+w))$ through order $v^4$ yields
$6(1+q)^dP/(4320q^4)$.

For $\zeta\ne1$, set $\alpha=1/(1-\zeta)$. Both $\frac vqa_\zeta$ and
$v(a_\zeta+b)-1$ expand as power series in $v$ whose coefficient of $v^k$ is a
polynomial in $\alpha$ of degree at most $k$:
\begin{align*}
 \frac vqa_\zeta&=\frac{\alpha}{q}v+\frac{\alpha-\alpha^2}{q^2}v^2
 +\frac{\alpha-3\alpha^2+2\alpha^3}{2q^3}v^3
 +\frac{\alpha-7\alpha^2+12\alpha^3-6\alpha^4}{6q^4}v^4+\cdots,\\
 v(a_\zeta+b)-1&=\Bigl(\alpha-\frac12\Bigr)v+\Bigl(\frac{\alpha-\alpha^2}{q}+\frac1{12}\Bigr)v^2
 +\frac{\alpha-3\alpha^2+2\alpha^3}{2q^2}v^3\\
 &\quad+\Bigl(\frac{\alpha-7\alpha^2+12\alpha^3-6\alpha^4}{6q^3}-\frac1{720}\Bigr)v^4+\cdots.
\end{align*}
Thus $[v^4]\frac vqa_\zeta(v(a_\zeta+b))^d$ is a polynomial in $\alpha$ of
degree at most $4$ with coefficients polynomial in $d$ and $1/q$. To sum this
polynomial over the nonprincipal roots, we compute the power sums
$p_k=\sum_{\zeta\ne1}\alpha^k$. Since the zeros of $((1+s)^q-1)/s$ are
$s=\zeta-1=-1/\alpha$, logarithmic differentiation gives
\[
 \frac{d}{ds}\log\frac{(1+s)^q-1}{s}=\sum_{\zeta\ne1}\frac{\alpha}{1+\alpha s}
 =\sum_{k\ge0}(-s)^kp_{k+1}.
\]
Extracting the first few coefficients gives $p_0=q-1$ and
\begin{gather*}
 p_1=\frac{q-1}2,\qquad p_2=-\frac{(q-1)(q-5)}{12},\qquad p_3=-\frac{(q-1)(q-3)}8,\\
 p_4=\frac{(q-1)(q^3+q^2-109q+251)}{720}.
\end{gather*}
Substituting these values of $p_k$ yields $(q-1)T/(4320q^4)$ for the sum over
$\zeta\ne1$, completing the proof.
\end{proof}

The algebraic expansion in this proof is routine. The script
\texttt{codim4\_symbolic.py} verifies it symbolically in $q$ and $d$, and
\texttt{rows.py} checks~\eqref{eq:codim4} against the $h^*$-polynomial for all
$5\le d\le20$ and $1\le q\le35$.

\begin{proof}[Proof of Theorem~\ref{thm:codim4}]
First assume $d\ge6$, and set $D=(5d-3)/2$. Write $X=6(1+q)^dP+(q-1)T$, which
has the same sign as $A_{d,4}$. Using $d+5\binom d2=d(5d-3)/2=dD$, we rewrite
\[
 P=-dq^2(q-D)+4dq-1.
\]
If $q\le D$, then $P\ge4dq-1\ge3dq$. If $q>D$, then $q-D\ge\frac12$ because
$2q$ and $5d-3$ are integers, and $q\ge14$ because $D\ge27/2$; hence
$P\le-\frac12dq^2+4dq-1\le-3dq$. In both cases $|P|\ge3dq$, and $P$ has the
sign claimed for $A_{d,4}$. It remains to prove that $|(q-1)T|<6(1+q)^d|P|$.

Replacing each coefficient of $T$ by its absolute value and using $d\ge1$ and
$q\ge1$ yields $|(q-1)T|\le219\,d^3q^7$.

\emph{Case $d\ge12$.} The function $q^6/(1+q)^d$ decreases for $q\ge1$ because
$6(1+q)\le dq$. Since $d^2/2^d$ decreases for $d\ge3$, we obtain
\[
 \frac{|(q-1)T|}{6(1+q)^d\cdot dq}\le\frac{219}{6}\,d^2\,\frac{q^6}{(1+q)^d}
 \le\frac{219}{6}\cdot\frac{d^2}{2^d}\le\frac{219}{6}\cdot\frac{144}{4096}<3.
\]
Therefore $|(q-1)T|<6(1+q)^d\cdot3dq\le6(1+q)^d|P|$.

\emph{Case $6\le d\le11$ and $q\ge100d$.} Here $q>D$ and
$q-D\ge q-\frac52d\ge\frac{39}{40}q$, so
$|P|\ge\frac{39}{40}dq^3-4dq\ge\frac12dq^3$. Using $(1+q)^d\ge q^6$, we get
\[
 \frac{|(q-1)T|}{6(1+q)^d}\le\frac{219}{6}d^3q<\frac12dq^3\le|P|,
\]
because $q^2\ge10^4d^2>73d^2$.

\emph{Case $6\le d\le11$ and $q<100d$.} Only finitely many pairs remain. For
each one, \texttt{rows.py} evaluates $X$ exactly and confirms its sign.

Finally, for $d\le5$, we list the coefficients explicitly. From~\eqref{eq:conv},
the rescaled coefficients $d!\,[t^k]L_d(t)$ for $0<k<d$, written as
polynomials in $r$, are
\begin{align*}
 d=2:&\quad r+2;\\
 d=3:&\quad \tfrac12(3r^2+r+18),\quad \tfrac32(r^2+r+2);\\
 d=4:&\quad 2(2r^2+r+22),\quad 4r^3+5r^2+2r+24,\quad 2(r^3+r^2+r+2);\\
 d=5:&\quad \tfrac16(-5r^4+60r^3+65r^2+24r+1500),\quad \tfrac{25}2(r^3+2r^2+r+14),\\
 &\quad \tfrac53(5r^4+6r^3+7r^2+3r+30),\quad \tfrac52(r^4+r^3+r^2+r+2).
\end{align*}
In each case the constant term is $1$ and the leading coefficient is $r^d/d!$.
Since every coefficient is a polynomial in $r$ of degree at most $d$, these
identities follow from their values at $d+1$ bases, which \texttt{rows.py}
checks. Every listed polynomial is positive except the linear coefficient for
$d=5$. By Descartes' rule of signs, $g(r)=-5r^4+60r^3+65r^2+24r+1500$ has
exactly one positive zero, and $g(13)>0>g(14)$.
\end{proof}

\section{Fixed codimension: proof of Theorem~\ref{thm:fixed}}\label{sec:fixed}

Write $U=\frac vqa_1$ and $G=v(a_1+b)/r$, so that the principal term of
Lemma~\ref{lem:residue} is $r^d[v^\ell]UG^d$. Setting $f(x)=x/(1-e^{-x})$, we
have $va_1=qf(v/q)$ and $f(x)=x/2+(x/2)\coth(x/2)$. Because $(x/2)\coth(x/2)$
is even, we obtain the coefficientwise expansions
\[
 U=\frac{v}{2q}+\Bigl(\text{even series}\Bigr),\qquad
 va_1+vb=q+\frac v2+\frac{v}{e^v-1}+O(q^{-1})=q+\frac v2\coth\frac v2+O(q^{-1}),
\]
where the error bound $O(q^{-1})$ holds uniformly for each fixed coefficient.
Moreover, $G$ is even because $a_1(-v)=1-a_1(v)$ and $b(-v)=-1-b(v)$. Its
constant term is $1$ and its coefficient of $v^2$ is $1/(12q)$.

\emph{Nonprincipal roots.} Let $\zeta\ne1$ and put $\alpha=1/(1-\zeta)$. Since
$|1-\zeta|\ge2\sin(\pi/q)\ge4/q$, we have $|\alpha|\le q/4$. Writing $s=v/q$
gives
\[
 a_\zeta=\frac{\alpha e^s}{1+\alpha(e^s-1)}=\alpha e^s\sum_{j\ge0}(-\alpha)^j(e^s-1)^j.
\]
Consequently, the coefficient of $v^k$ in $\frac vqa_\zeta$ is a polynomial in
$\alpha$ of degree at most $k$ divided by $q^k$, hence $O_k(1)$. Similarly,
the coefficient of $v^k$ in $v(a_\zeta+b)-1$ is $O_k(q)$. Expanding the $d$th
power shows that each nonprincipal term of Lemma~\ref{lem:residue} has
absolute value at most $C_\ell(qd)^\ell$. Summing over all nonprincipal roots
bounds their total by $C_\ell\,q\,(qd)^\ell$, where $C_\ell$ depends only on
$\ell$.

\emph{Part (1).} Suppose $r/d\to c$. Then the nonprincipal contribution
to $r^{-d}A_{d,\ell}$ is $O(r^{-d}q^{\ell+1}d^\ell)$, which tends to $0$.
For the principal term, we have $U\to1$ and
\[
 d\log G=\frac dr\Bigl(\frac v2\coth\frac v2-1\Bigr)+O\Bigl(\frac d{r^2}\Bigr)
 \longrightarrow\frac1c\Bigl(\frac v2\coth\frac v2-1\Bigr)
\]
coefficientwise, which establishes the first limit. Because the limit is even
in $v$, all of its odd coefficients vanish. For odd $\ell$, the parity of $G$
and $U-v/(2q)$ yields $[v^\ell]UG^d=\frac1{2q}[v^{\ell-1}]G^d$, and we have
$r/(2q)\to\frac12$. Meanwhile, the nonprincipal contribution to
$r^{1-d}A_{d,\ell}$ is still $O(r^{1-d}q^{\ell+1}d^\ell)\to0$, giving the
second limit.

\emph{Part (2).} Now fix $r$. When $\ell=2j$, the coefficient $[v^{2j}]G^d$ is
a polynomial in $d$ of degree $j$ with leading term $\binom dj(12q)^{-j}$,
because achieving degree $j$ in $d$ requires choosing the $v^2$ term from $j$
separate factors. Since $U$ has constant term $1$, $[v^\ell]UG^d$ is a
polynomial in $d$ with positive leading coefficient. When $\ell=2j+1$,
$[v^\ell]UG^d=\frac1{2q}[v^{2j}]G^d$ has the same property. Thus the principal
term equals $r^d$ times a polynomial in $d$ with positive leading coefficient,
whereas the nonprincipal terms are bounded by a polynomial in $d$ since $q$ is
fixed. Hence $A_{d,\ell}>0$ for all sufficiently large $d$. \qed

\section{Bases two and three}\label{sec:23}

To prove Ehrhart positivity for small bases, we bundle all dimensions together
into a single generating function. Summing~\eqref{eq:gf} over $d$ with a new
variable $z$ gives
\begin{equation}\label{eq:sumd}
 \sum_{d\ge0}L_d(t)z^d=[x^{qt}]\,\frac{1}{1-x}\cdot\frac{1}{1-zS_q(x)/(1-x^q)}
 =[x^{qt}]\,\frac{S_{q-1}(x)}{1-x^q-zS_q(x)},
\end{equation}
as formal power series in $z$. At each power $z^d$, both sides are polynomials
in $t$ that agree for all integers $t\ge0$.

\begin{theorem}\label{thm:23}
For $r=2$ and $r=3$, every coefficient of $L_d(t)$ is positive for every $d\ge1$.
\end{theorem}

\begin{proof}
\emph{The case $q=1$.} Here~\eqref{eq:sumd} reduces to
$[x^t]\,1/((1-z)-x(1+z))$, so
\[
 \sum_{d\ge0}L_d(t)z^d=\frac1{1-z}\Bigl(\frac{1+z}{1-z}\Bigr)^t
 =\frac1{1-z}\exp\bigl(t\lambda(z)\bigr),\qquad
 \lambda(z)=\log\frac{1+z}{1-z}=2\sum_{m\ge0}\frac{z^{2m+1}}{2m+1}.
\]
Extracting coefficients gives $k!\,[t^kz^d]=[z^d](1-z)^{-1}\lambda(z)^k$. The
series $\lambda^k$ has nonnegative coefficients with $2^k$ at $z^k$, while
every coefficient of $(1-z)^{-1}$ is $1$. Hence the coefficient is positive
for all $0\le k\le d$.

\emph{The case $q=2$.} Now~\eqref{eq:sumd} becomes
$[x^{2t}]\,(1+x)/(1-x^2-z(1+x+x^2))$. The denominator factors as
$-(1+z)(x-x_+)(x-x_-)$, where $x_\pm=(-z\pm\sqrt{4-3z^2})/(2(1+z))$ satisfies
$x_\pm(0)=\pm1$. Expanding in partial fractions gives
\[
 \sum_{d\ge0}L_d(t)z^d=b_+(z)e^{t\lambda_+(z)}+b_-(z)e^{t\lambda_-(z)},\qquad
 \lambda_\pm=-2\log(\pm x_\pm),
\]
where, after simplification,
\[
 b_\pm=\frac{1}{2(1-z)}\Bigl(1\pm\frac{2-z}{\sqrt{4-3z^2}}\Bigr),\qquad
 \lambda_\pm'=\frac{2}{1-z^2}\Bigl(1\pm\frac1{\sqrt{4-3z^2}}\Bigr),\qquad \lambda_\pm(0)=0.
\]
As in the case $q=1$, this identity holds at each power of $z$ for all
integers $t\ge0$, and therefore holds as an equality of polynomials in $t$.

Set $a_i=[z^{2i}](4-3z^2)^{-1/2}=\frac12\binom{2i}{i}(3/16)^i>0$. Evaluating
at $z=1$ gives $\sum_ia_i=1$, and the consecutive ratios satisfy
$a_{i+1}/a_i=\frac{3(2i+1)}{8(i+1)}$, which is at least $\frac9{16}$ for all
$i\ge1$. Multiplying either $1+(4-3z^2)^{-1/2}$ or
$1-(4-3z^2)^{-1/2}=\frac12-\sum_{i\ge1}a_iz^{2i}$ by $(1-z^2)^{-1}$ produces a
series with nonnegative coefficients: in the second product, the coefficient
of $z^{2m}$ is $1-\sum_{i\le m}a_i\ge0$. Integrating term by term shows that
both $\lambda_+$ and $\lambda_-$ have nonnegative coefficients, with
$[z]\lambda_+=3$.

To analyze $b_+$, rewrite $\frac{2-z}{1-z}=1+\frac1{1-z}$. Then
$2b_+=\frac1{1-z}+\bigl(1+\frac1{1-z}\bigr)(4-3z^2)^{-1/2}$ clearly has
positive coefficients. For $b_-$, the series
$2(1-z)b_-=1-(2-z)(4-3z^2)^{-1/2}$ has coefficients $1-2a_0=0$ at $z^0$, $a_m$
at each odd degree $2m+1$, and $-2a_m$ at each even degree $2m\ge2$. Dividing
by $1-z$ replaces these coefficients by their partial sums. Through degree
$2m+1$, the partial sum is
$1-2\sum_{i\le m}a_i+\sum_{i\le m}a_i=1-\sum_{i\le m}a_i\ge0$. Through degree
$2m$, it is $1-\sum_{i\le m}a_i-a_m=\sum_{i>m}a_i-a_m$, which equals $0$ when
$m=0$. For $m\ge1$, the ratio bound gives
$\sum_{i>m}a_i\ge a_m\sum_{j\ge1}(9/16)^j=\frac97a_m>a_m$. Thus $b_-$ has
nonnegative coefficients.

Finally, extracting $k!\,[t^kz^d]$ from the sum gives
$[z^d]\,b_+\lambda_+^k+[z^d]\,b_-\lambda_-^k$. The second term is nonnegative.
In the first term, $\lambda_+^k$ has nonnegative coefficients with $3^k$ at
$z^k$, while $b_+$ has positive coefficients. Their product is therefore
positive for all $k\le d$.
\end{proof}

\section{Bases four and five}\label{sec:45}

\begin{theorem}\label{thm:45}
For $r=4$ and $r=5$, every coefficient of $L_d(t)$ is positive for every $d\ge1$.
\end{theorem}

Throughout this section, let $q\in\{3,4\}$.

\subsection{Branches}
For each $q$th root of unity $\zeta$, let $x_\zeta(z)$ be the root of
$1-x^q-zS_q(x)=0$ with $x_\zeta(0)=\zeta$. Because these roots are simple at
$z=0$, each $x_\zeta(z)$ expands as a power series in $z$. Expanding~\eqref{eq:sumd}
in partial fractions and using $[x^n](x-x_\zeta)^{-1}=-x_\zeta^{-n-1}$ gives
\begin{equation}\label{eq:branches}
 \sum_{d\ge0}L_d(t)z^d=\sum_{\zeta^q=1}b_\zeta(z)e^{tv_\zeta(z)},\qquad
 v_\zeta=-q\log\frac{x_\zeta}{\zeta},\qquad
 b_\zeta=-\frac{z\,x_\zeta'}{x_\zeta(1-x_\zeta)}.
\end{equation}
To obtain the formula for $b_\zeta$, we rewrite the defining equation as
$S_{q-1}(x_\zeta)=zS_q(x_\zeta)/(1-x_\zeta)$ and differentiate implicitly to
get $\partial_x(1-x^q-zS_q)\cdot x_\zeta'=S_q(x_\zeta)$. As before, this
identity holds at each power of $z$ as an equality of polynomials in $t$. From
$x_1=1-\frac{q+1}{q}z+O(z^2)$, we find $b_1(0)=1$ and $[z]v_1=q+1$, whereas
$b_\zeta(0)=v_\zeta(0)=0$ for every $\zeta\ne1$. Consequently,
\begin{equation}\label{eq:tk}
 k!\,[t^kz^d]\sum_{d}L_d(t)z^d=\sum_{\zeta^q=1}[z^d]\,b_\zeta v_\zeta^k .
\end{equation}
Our strategy is to show that the principal term $\zeta=1$ dominates the sum.

Notice also that replacing $x$ by $1/x$ and $z$ by $-z$ preserves the root
equation, and the root equal to $1$ at $z=0$ is simple, so $x_1(-z)=1/x_1(z)$ and
$v_1$ is odd.

\subsection{Coefficients by Lagrange inversion}
We first compute the initial coefficients of every branch exactly via Lagrange
inversion. Write $x_\zeta=\zeta(1-y)$. The root equation then takes the form
$y=z\phi(y)$, where
\[
 \phi(y)=S_q\bigl(\zeta(1-y)\bigr)\cdot\frac{y}{1-(1-y)^q}.
\]
Since $zy'=y/(1-z\phi'(y))$, Lagrange inversion yields
$[z^n]\,zy'\,H(y)=[y^{n-1}]H(y)\phi(y)^n$ for every power series $H$ and every
$n\ge1$. In this coordinate we have $v_\zeta=-q\log(1-y)$,
$b_\zeta=zy'/((1-y)(1-\zeta+\zeta y))$ for $\zeta\ne1$, and, because
$y/z=\phi(y)$, $b_1=1+z\frac{d}{dz}\log\frac{\phi(y)}{1-y}$. Applying Lagrange
inversion gives, for all $n\ge1$,
\begin{equation}\label{eq:lagrange}
 [z^n]v_\zeta=\frac qn[y^{n-1}]\frac{\phi^n}{1-y},\qquad
 [z^n]b_1=[y^n]\frac{\phi^n}{1-y},\qquad
 [z^n]b_\zeta=[y^{n-1}]\frac{\phi^n}{(1-y)(1-\zeta+\zeta y)}\ (\zeta\ne1).
\end{equation}
In the expression for $b_1$, the term $\phi'/\phi$ contributes
$[y^{n-1}]\phi'\phi^{n-1}=[y^n]\phi^n$.

To carry out these computations exactly, substitute $y=qu$ and set
$A(u)=S_q(\zeta(1-qu))$, $B(u)=(1-(1-qu)^q)/(qu)$, and $F=(qA/B)^n=q^n\phi^n$.
The coefficients of $F$ are determined by $F(0)=A(0)^n$ and the linear
recurrence encoded in
\[
 A\,B\,F'=n\,(A'B-AB')\,F .
\]
Writing $C_j=[u^j]F/(1-qu)$, the formulas in~\eqref{eq:lagrange} become
\[
 [z^n]v_\zeta=\frac{q^2C_{n-1}}{nq^{2n}},\qquad [z^n]b_1=\frac{C_n}{q^{2n}},\qquad
 [z^n]b_\zeta=\frac{\beta D_{n-1}}{q^{2n}}\ (\zeta\ne1),
\]
where $\beta=q/(1-\zeta)$, $D_{-1}=0$, and $D_j=C_j-\zeta\beta D_{j-1}$. Every
quantity in this recurrence lies in the ring $\Z[\zeta]$ (the Eisenstein
integers, the Gaussian integers, or $\Z$). The script \texttt{bases\_4\_5.py}
runs the recurrence exactly for one branch from each complex-conjugate pair,
verifying divisibility at every step. It also checks that~\eqref{eq:branches}
reproduces $L_d(t)$ for all $d\le12$.

\begin{lemma}[finite certificate]\label{lem:finite}
Let $N=180$ for $q=3$ and $N=900$ for $q=4$. For $1\le n\le N$:
\begin{enumerate}
\item $[z^n]b_1\ge\frac12$, and $[z^n]v_1>0$ for odd $n$;
\item $|[z^n]b_\zeta|<\frac15$ for $\zeta\ne1$;
\item $\sum_{n=1}^N|[z^n]v_\zeta|$ is less than $\frac74$ for $q=3$, $\zeta\ne1$;
less than $\frac{12}5$ for $q=4$, $\zeta=\pm i$; and less than $\frac54$ for
$q=4$, $\zeta=-1$.
\end{enumerate}
\end{lemma}

\begin{proof}
These bounds follow from an exact computation using the formulas above. Upper
bounds on absolute values are obtained via the norm form of $\Z[\zeta]$ and
integer square roots, rounded up to multiples of $10^{-9}$. The computed
values are $\min_n[z^n]b_1\approx0.667$ and $0.625$,
$\max|[z^n]b_\zeta|\approx0.193$ ($q=3$) and $0.177$, $0.125$ ($q=4$,
$\zeta=i,-1$), while the sums in~(3) are approximately $1.716$, $2.356$, and
$1.210$.
\end{proof}

\subsection{Analytic continuation}
To bound the remaining coefficients, we analytically continue the branches to
a closed disk $|z|\le R$, with $R=9/8$ for $q=3$ and $R=51/50$ for $q=4$.
Setting $x=(1-w)/(1+w)$ transforms the inverse relation $z=(1-x^q)/S_q(x)$
into
\[
 z=\frac{w(3+w^2)}{2(1+w^2)}\quad(q=3),\qquad
 z=\frac{8w(1+w^2)}{5+10w^2+w^4}\quad(q=4).
\]
Differentiating shows that the finite critical points satisfy $w^4+3=0$ for
$q=3$, and $u^3-7u^2-5u-5=0$ with $u=w^2$ for $q=4$.

For $q=3$, setting $w^2=\pm i\sqrt3$ gives $|3+w^2|^2=12$ and $|1+w^2|^2=4$.
Every critical value therefore has modulus $3^{3/4}/2$, which lies outside the
disk since $(3^{3/4}/2)^4=27/16>(9/8)^4$.

For $q=4$, the cubic has a unique real zero $U\in(7.73045,7.73047)$; it has no
zero $u\le0$ because $u^3\le0<7u^2+5u+5$ there. The other two zeros form a
complex-conjugate pair $a\pm ib$ satisfying $2a+U=7$ and $(a^2+b^2)U=5$, which
forces $a\in(-0.37,-0.36)$. The squared modulus of a critical value is
$64|u|\,|1+u|^2/|5+10u+u^2|^2$. At $u=U$, this exceeds $1.36^2$. At
$u=a\pm ib$, setting $m=a^2+b^2$ gives
\[
 \frac{64\sqrt m\,(1+2a+m)}{(5+10a+2a^2-m)^2+(10+2a)^2(m-a^2)}>\frac{47.15}{45.06}>\Bigl(\frac{51}{50}\Bigr)^2,
\]
by interval bounds derived from the bracket for $U$ (script \texttt{disks.py}).

The rational map is also locally injective at its poles, which are simple, and
at $w=\infty$: in the local coordinate $1/w$, the point $w=\infty$ is a simple
pole for $q=3$ and a simple zero for $q=4$. Thus no critical value lies in the
closed disk $|z|\le R$, and every branch $x_\zeta$ extends meromorphically to
a neighborhood of that disk.

We now track where the branches take the values $0$, $1$, and~$\infty$. A root
$x$ of $1-x^q-zS_q(x)$ can equal $0$ only when $z=1$, can equal $\infty$ only
when $z=-1$ (where the leading coefficient $-(1+z)$ vanishes), and can equal
$1$ only when $z=0$. The value $0$ is attained by $x_1$ at $z=1$: along
$w\in[0,1]$, the map increases from $0$ to $1$ without critical points while
$x$ moves from $1$ to $0$. By the symmetry $x_1(-z)=1/x_1(z)$, the pole at
$z=-1$ also lies on $x_1$, and both singularities are simple. Consequently, on
the disk, every nonprincipal branch $x_\zeta$ for $\zeta\ne1$ is holomorphic
and avoids $0$, $1$, and~$\infty$.

\subsection{Tail bounds}
To isolate the singularities of the principal branch on $|z|\le R$, we define
holomorphic logarithms vanishing at $z=0$ by
\begin{align*}
 V_1&=-q\log\Bigl(x_1\frac{1+z}{1-z}\Bigr), &
 W_1&=\log\Bigl(\frac{x_1}{1-x_1}\cdot\frac{z}{1-z}\cdot\frac{q+1}{q}\Bigr),\\
 V_\zeta&=v_\zeta, &
 W_\zeta&=\log\Bigl(\frac{x_\zeta}{1-x_\zeta}\Big/\frac{\zeta}{1-\zeta}\Bigr)\quad(\zeta\ne1).
\end{align*}
Every argument inside these logarithms is holomorphic and zero-free on the
disk: the zero of $x_1$ at $z=1$ cancels against $1-z$, the pole at $z=-1$
cancels against $1+z$ or $1-x_1$, and the zero of $1-x_1$ at $z=0$ cancels
against $z$. In terms of these functions, we have
\begin{equation}\label{eq:VW}
 v_1=q\log\frac{1+z}{1-z}+V_1,\qquad b_1=\frac1{1-z}-zW_1',\qquad b_\zeta=-zW_\zeta'\ (\zeta\ne1).
\end{equation}

\begin{lemma}\label{lem:tail}
For $n\ge1$, $|[z^n]V_\zeta|\le100R^{-n}$ and $|[z^n]W_\zeta|\le100R^{-n}$ for every $\zeta$.
\end{lemma}

\begin{proof}
Suppose $|z|=R$. Dividing $1-x^q-zS_q(x)$ by its leading coefficient $-(1+z)$
yields a monic polynomial whose remaining coefficients have modulus at most
$\max(|z|,|1-z|)/|1+z|\le(1+R)/(R-1)\le101$. The same bound holds for the
reversed polynomial after dividing by $1-z$. Cauchy's root bound therefore
gives $1/150\le|x_\zeta|\le150$. Evaluating at $x=1$ gives
$\prod_\zeta(1-x_\zeta)=(q+1)z/(1+z)$, so
$\prod_\zeta|1-x_\zeta|\ge(q+1)R/(1+R)\ge1$. Combining this with
$|1-x_\zeta|\le151$ yields $|1-x_\zeta|\ge151^{-3}>10^{-7}$.

In addition, $1/50\le|1\pm z|\le17/8$, and $\sqrt2\le|1-\zeta|\le2$ for all
$\zeta\ne1$. Taking logarithms of the moduli gives
$|\operatorname{Re}V_\zeta|\le50$ and $|\operatorname{Re}W_\zeta|\le50$ on the
circle. Now, for any function $f$ holomorphic on $|z|\le R$ with $f(0)=0$, the
$n$th Fourier coefficient of $\operatorname{Re}f(Re^{i\theta})$ is
$\frac12R^n[z^n]f$. Hence
$|[z^n]f|\le2R^{-n}\max|\operatorname{Re}f|\le100R^{-n}$.
\end{proof}

\begin{proof}[Proof of Theorem~\ref{thm:45}]
For all $n>N$, we have $100nR^{-n}<\frac1{100}$ and
$\sum_{n>N}100R^{-n}<\frac1{100}$. Both inequalities are verified in exact
rational arithmetic at $n=N+1$, and $100nR^{-n}$ decreases for $n\ge N$.
Applying~\eqref{eq:VW} and Lemma~\ref{lem:tail} for $n>N$ gives four bounds:
$[z^n]b_1\ge1-100nR^{-n}>\frac12$; $[z^n]v_1\ge\frac{2q}n-100R^{-n}>0$ for odd
$n$; $[z^n]v_1=0$ for even $n$ since $v_1$ is odd; and
$|[z^n]b_\zeta|\le100nR^{-n}<\frac1{100}$ for $\zeta\ne1$. Combining these
tail estimates with Lemma~\ref{lem:finite}, we conclude that every coefficient
of $b_1$ is at least $\frac12$, every coefficient of $v_1$ is nonnegative,
$|[z^n]b_\zeta|<\frac14$ for all $n$ and $\zeta\ne1$, and
\[
 \sum_{n\ge1}|[z^n]v_\zeta|<\begin{cases}2,&q=3,\\ \frac52,&q=4.\end{cases}
\]

Now let $1\le k\le d$. Because $b_1$ and $v_1$ have nonnegative coefficients
and $[z]v_1=q+1$, the principal term satisfies
\[
 [z^d]\,b_1v_1^k\ge[z^{d-k}]b_1\cdot(q+1)^k\ge\tfrac12(q+1)^k.
\]
For each nonprincipal root $\zeta\ne1$, we have $v_\zeta(0)=0$, so
$|[z^d]b_\zeta v_\zeta^k|\le\frac14\bigl(\sum_n|[z^n]v_\zeta|\bigr)^k$.
Summing over the $q-1$ nonprincipal roots bounds their total below
$\frac12\cdot2^k$ for $q=3$ and below $\frac34(\frac52)^k$ for $q=4$. Both
bounds are smaller than $\frac12(q+1)^k$, which equals $\frac12\cdot4^k$ or
$\frac12\cdot5^k$. By~\eqref{eq:tk}, it follows that $[t^k]L_d(t)>0$. Finally,
the constant term is $[t^0]L_d=L_d(0)=1$.
\end{proof}

\section{Proof of Theorem~\ref{thm:classification}}\label{sec:classification}

Theorems~\ref{thm:23} and~\ref{thm:45} establish Ehrhart positivity in every
dimension for $r=2,3,4,5$. By Theorem~\ref{thm:codim4}, every simplex
$B_{r,d}$ with $d\le4$ is Ehrhart positive, and for $r\ge14$ the linear
coefficient of $L_5$ is negative while all remaining coefficients are
positive. For $6\le r\le13$, the script \texttt{rows.py} computes $L_d(t)$
exactly from the $h^*$-polynomial for every $d$ up to the entry listed in
Table~\ref{tab:first}. In each case, all coefficients are positive below that
dimension, and at that dimension the linear coefficient is the sole negative
coefficient. \qed

Note that Table~\ref{tab:first} records only the first failure for each base,
not the set of all dimensions in which Ehrhart positivity fails.

\section*{Verification code}

The scripts named above are available in the repository
\begin{center}\url{https://github.com/brandonlign/base-r-ehrhart}.\end{center}
All of them except
\texttt{codim4\_symbolic.py} use only the Python standard library and exact
integer or rational arithmetic; \texttt{run\_all.sh} runs them in about
fifteen seconds.
\begin{itemize}
\item \texttt{ehrhart.py}: computes $L_d(t)$ from the
$h^*$-polynomial~\cite{Sol}, from the convolution~\eqref{eq:conv}, and by
direct lattice-point counts, comparing them on 447 values, including Graham's
$L_5$ for $r=14$.
\item \texttt{rows.py}: verifies the signs in Theorem~\ref{thm:large} for
$d\le16$, the formula~\eqref{eq:codim4} for $5\le d\le20$ and $q\le35$, the
finite cases of Theorem~\ref{thm:codim4}, the formulas for $d\le5$,
Table~\ref{tab:first}, and the coefficient $(5-2c)/(1440c^2)$.
\item \texttt{bases\_2\_3.py}: checks the generating functions of
Section~\ref{sec:23} for $d\le14$ and the coefficient bounds used there
through degree $200$.
\item \texttt{bases\_4\_5.py}: verifies Lemma~\ref{lem:finite}, the tail
inequalities, and the identity~\eqref{eq:branches} for $d\le12$.
\item \texttt{disks.py}: checks the maps in $w$ and the critical-value bounds.
\item \texttt{codim4\_symbolic.py} (needs \texttt{sympy}, runs in a few
minutes): carries out the expansion in the proof of
Proposition~\ref{prop:codim4} symbolically in $q$ and $d$.
\end{itemize}

\section*{Acknowledgments}
Claude (Anthropic) helped with the proofs, the
verification code, and the manuscript.

\begin{thebibliography}{9}

\bibitem{BCKV}
D.~Bump, K.-K.~Choi, P.~Kurlberg and J.~Vaaler, \emph{A local Riemann
hypothesis, I}, Math. Z. \textbf{233} (2000), 1--19.

\bibitem{Gra}
P.~Graham, \emph{A base-$r$ simplex that is not Ehrhart positive},
Hexagon 2611.00793v1 (21 November 2026),
\url{https://hexagonmath.org/2611.00793v1};
\href{https://doi.org/10.5281/zenodo.23875145}{doi:10.5281/zenodo.23875145}.

\bibitem{Liu}
F.~Liu, \emph{On positivity of Ehrhart polynomials}, in: Recent Trends in
Algebraic Combinatorics, Assoc. Women Math. Ser.~16, Springer, 2019,
189--237; \href{https://arxiv.org/abs/1711.09962}{arXiv:1711.09962}.

\bibitem{LS}
F.~Liu and L.~Solus, \emph{On the relationship between Ehrhart unimodality
and Ehrhart positivity}, Ann. Comb. \textbf{23} (2019), 347--365;
\href{https://arxiv.org/abs/1804.08258}{arXiv:1804.08258}.

\bibitem{Sol}
L.~Solus, \emph{Simplices for numeral systems}, Trans. Amer. Math. Soc.
\textbf{371} (2019), 2089--2107;
\href{https://arxiv.org/abs/1706.00480}{arXiv:1706.00480}.

\end{thebibliography}

\end{document}
